\documentclass{article}
\usepackage[T1]{fontenc}
\usepackage{lmodern}
\usepackage{graphicx}
\usepackage{amsmath,amssymb}
\usepackage[a4paper,margin=2cm]{geometry}
\usepackage[absolute,overlay]{textpos}
\setlength{\TPHorizModule}{1mm}
\setlength{\TPVertModule}{1mm}
\usepackage[indonesian]{babel}
\usepackage{hyperref}
\setlength{\emergencystretch}{2em}
% unit: Halaman 59

\begin{document}

% === Halaman 59 (llm) ===

\section*{4. Enkripsi}

\begin{itemize}
    \item Pesan dapat dienkripsi terlebih dahulu sebelum disembunyikan ke dalam citra.
    \item Teknik enkripsi yang sederhana misalnya dengan meng-XOR-kan bit-bit pesan dengan bit-bit kunci. Jumlah bit-bit kunci sama dengan jumlah bit pesan.
    \item Bit-bit kunci dibangkitkan secara acak.
    \item Kunci untuk pembangkitan bit-bit kunci menjadi \textit{stego-key}.
    \item Jika dipakai teknik acak dalam memilih \textit{pixel-pixel}, maka ada dua \textit{stego-key}: satu untuk pembangkitan bit-bit kunci, satu lagi untuk pembangkitan posisi \textit{pixel} yang dipilih untuk menyembunyikan pesan.
\end{itemize}

\end{document}
